\documentclass[dvipsnames,12pt]{beamer}
\usepackage{amsmath,amssymb}
\usepackage{amsfonts}
\usepackage{colortbl}
\usepackage[english]{babel}
\usepackage[brazil]{varioref}
\usepackage{xspace}
\usepackage{cmap}				% Mapear caracteres especiais no PDF
\usepackage{lmodern}
\usepackage[utf8]{inputenc}
\usepackage{epstopdf}
\usepackage{ragged2e}
\usepackage{graphics}
\usepackage{wasysym}
\usepackage{ucs}
\usepackage{pgf}
\usepackage{verbatim}
\usepackage{enumerate}
\usepackage{float, subfigure, multicol, xspace, fancyhdr, wasysym,url}
\usepackage[export]{adjustbox}  % Ajusta o alinhamento entre as figuras
% % TM MARK
% \usepackage{textcomp}
\usepackage{multirow}
% PGFPlot packages
\usepackage{tikz} % Pacotes para plotar gráficos a partir de arquivos de texto
\usepackage{transparent}
\usepackage{mathpazo}

\usepackage{pgfplots}
\usepackage{pgfplotstable}
\usepackage{siunitx}
% \usepgfplotslibrary{external}
%\tikzexternalize[prefix=TikzPictures/] % activate externalization!
\pgfplotsset{compat=newest} % Allows to place the legend below plot
\usepgfplotslibrary{units} % Allows to enter the units nicely
\sisetup{
  round-mode          = places,
  round-precision     = 2,
}
\usepackage{tabularx}
\newcolumntype{C}{>{\centering\arraybackslash}X}
\renewcommand{\arraystretch}{1.2}
\usepackage{circuitikz}
\newcommand{\derivOne}[2]{\dfrac{\partial #1}{\partial #2}}

\mode<presentation> {\usetheme{LVA}}
% \mode<presentation> {\usetheme{CambridgeUS}}
\usecolortheme{orchid}
% Tira os símbolos de navegação
\setbeamertemplate{navigation symbols}{}
% Enumera as figuras
\setbeamertemplate{caption}[numbered]

\makeatletter
\newenvironment{noheadline}{
    \setbeamertemplate{headline}{}
    \addtobeamertemplate{frametitle}{\vspace*{-0.9\baselineskip}}{}
}{}
\makeatother



% \usetikzlibrary{shapes.geometric, arrows}
% \tikzstyle{startstop} = [rectangle, rounded corners, minimum width=3cm, minimum height=1cm,text centered, draw=black, fill=darkblue!30]
% \tikzstyle{arrow} = [thick,->,>=stealth]
%\setbeamertemplate{caption}[numbered]

% \usefonttheme[onlymath]{serif}

\providecommand{\sin}{} \renewcommand{\sin}{\hspace{2pt}\mathrm{sen}}		% Redefine a escrita do seno
\providecommand{\tan}{} \renewcommand{\tan}{\hspace{2pt}\mathrm{tg}}		% Redefine a escrita da tangente


%%%%%%%%%%%%%%%%%%%%%%%%%%%%%%%%%%%%%%%%%%%%%%%%%%%%%%%%%%%%%%%%%%%%%
\title[Laboratory of Acoustic and Vibrations]{Numerical investigation of normal mode radiation properties of ducts with low Mach number inlet flow}
\author[José Pedro de Santana Neto]{José Pedro de Santana Neto*; Danilo Braga, Andrey Ricardo da Silva; Julio Cordioli}
\institute[LVA-UFSC]{Laboratory of Acoustic and Vibrations - Federal University of Santa Catarina}
\date{}
% \date{Buenos Aires, September, 2016}
 %\titlegraphic{\includegraphics[width=0.8\linewidth,valign=m]{./Figures/ICA_footer.pdf}}
%               \includegraphics[scale=0.45,valign=m]{./Figures/otovidalogo}~\hfill
% 			  \includegraphics[scale=0.20,valign=m]{./Figures/LVA_Logo.pdf}~\hfill
% 			  \includegraphics[scale=0.27,valign=m]{./Figures/wavetechlogo}~\hfill
% 			  \includegraphics[scale=0.30,valign=m]{./Figures/finep}}
%%%%%%%%%%%%%%%%%%%%%%%%%%%%%%%%%%%%%%%%%%%%%%%%%%%%%%%%%%%%%%%%%%%%%

%% Lista de novos comandos
\newcommand\etal{\textit{et al.}\xspace}

\newcommand{\fig}[1]{Figura~\ref{#1}\xspace}
\newcommand{\figs}[2]{Figuras~\ref{#1}~e~\ref{#2}\xspace}
%%%%%%%%%%%%%%%
%%%%%%%%%%%%%%%%
\newcommand{\tab}[1]{Tabela~\ref{#1}\xspace}
\newcommand{\tabs}[1]{Tabelas~\ref{#1}\xspace}
%%%%%%%%%%%%%%%
%%%%%%%%%%%%%%%%
% % % % % % % % % % % % % % % % % % % % % % % % % % % % % % % % % % % % % % % % % %

\begin{document}
\section*{}
\begin{noheadline}

\setbeamertemplate{footline}{} 

\begin{frame}
  
   \titlepage
   
\end{frame}

\addtocounter{framenumber}{-1}

\end{noheadline}

% \begin{frame}
%   \frametitle{Sumário}
%   \tableofcontents[section=1,hidesubsections]
% \end{frame}


\AtBeginSection[]
{
  \frame<handout:0>
  {
    \frametitle{Summary}
    \setbeamertemplate{footline}{} 
    \tableofcontents[currentsection,hideallsubsections]
  }
}

\AtBeginSubsection[]
{
  \frame<handout:0>
  {
    \frametitle{Sumário}
    \tableofcontents[sectionstyle=show/hide,subsectionstyle=show/shaded/hide]
  }
}

\newcommand<>{\highlighton}[1]{%
  \alt#2{\structure{#1}}{{#1}}
}
\newcommand{\icon}[1]{\pgfimage[height=1em]{#1}}

% %%%%%%%%%%%%%%%%%%%%%%%%%%%%%%%%%%%%%%%%%%%%%%%%%%%%%%%%%%%%%%%%%%%%%

\section{Introduction}
\begin{frame}
\frametitle{Context}
\begin{itemize}
    \item Ducts are presents in a several acoustics components nowadays; 
    \item Most of ducts have a inner mean flow;
    \item Study the acoustics properties of ducts with flow are important and there are so many works with outlet flow;
    \item But few works have presented concepts and informations internal acoustic about ducts with inlet flow;
    \item The numerical aproach by using the lattice Boltzmann method is a good choice to explore the acoustics parameters of dutcs with flow.  
\end{itemize}
\end{frame}

\begin{frame}
\frametitle{Goals}
\begin{itemize}
    \item Main goal: investigate the internal acoustic of unflanged ducts with normal mode and low Mach ($M \leq 0.2$) number inlet flow by using the lattice Boltzmann method. 
    \item Specific goals:
    \begin{itemize}
      \item implement and validate the internal acoustics parameters of unflanged duct without flow;
      \item implement and validate the internal acoustics parameters of unflanged duct with outlet flow;
      \item implement, validate and analyse the internal acoustics parameters of unflanged duct with inlet flow.
    \end{itemize}
\end{itemize}
\end{frame}


\section{Background}
\begin{frame}
\frametitle{Internal Acoustics Parameters}
\begin{itemize}
    \item Magnitude of coeficient reflection in end of duct:
     \begin{equation}
        |R_{r}| =\bigg|\frac{Z_{r} - Z_{0}}{Z_{r} + Z_{0}}\bigg|,
        \label{eq:R}
    \end{equation}
    \item End correction:
    \begin{equation}
        l = \frac{1}{k} \arctan\!\left(\frac{Z_{r}}{Z_{0} \, \mathrm{j}}\right).
        \label{eq:l}
    \end{equation}
\end{itemize}
\end{frame}

\begin{frame}
\frametitle{Related Works}
\begin{itemize}
    \item Levine and Schwinger (1948): $|R_{r}|$ and $l/a$ without flow;
    \item Munt (1990): $|R_{r}|$ and $l/a$ with outlet flow;
    \item Ingard and Singhal (1975) and Davies (1987): $|R_{r}|$ and $l/a$ with inlet flow but only low frequencies ($ka < 0.25 $). The values of $l/a$ was suggested but not proved or measured;
    \item Ingard and Singhal (1975) and Davies (1987) related that inlet flow form a vena contracta with factor loss defined by
    \begin{equation}
      K_{p} = \frac{\Delta p}{0.5\rho (c_{0}M)^{2}}.
    \end{equation} 
\end{itemize}
\end{frame}

\section{Lattice Boltzmann Method}
\begin{frame}
\frametitle{Governing Equations}
\begin{itemize}
    \item lattice Boltzmann with MRT model according d’Humieres (1994):
    \begin{equation}
      f_{i}(\textbf{x} + c_{i}\Delta t, t + \Delta t) = f_{i}(\textbf{x}, t) -\textbf{$M$}^{-1}\textbf{$S$}(m_{i} - m_{i}^{M}).
    \label{eq:MRT_5}
    \end{equation}
    \item rigid wall: bounceback no-slip Viggen (2014);
    \item anechoic condition: Kam et al. (2006). 
\end{itemize}
\end{frame}

\begin{frame}
\frametitle{Element Structure}
\begin{itemize}
    \item D3Q19:
    \begin{figure}[ht!]
    \centering
  \includegraphics[width=.75\linewidth]{figuras/d3q19.pdf}
  \caption[Esquemático do D3Q19]{D3Q19 Model. Figure adapted from Premnath et al. (2013).}
  \label{fig:d3q19}
\end{figure} 
\end{itemize}
\end{frame}

\begin{frame}
\frametitle{Numerical Model}
\begin{figure}[ht!]
\centering
  \includegraphics[width=1.\linewidth]{figuras/modelo_numerico_2.pdf}
  \caption{Numerical model: view of 3D lateral cut of the  model.}
  \label{fig:modelo}
\end{figure}
\end{frame}

\newcommand\scalex{1}
\newcommand\scaley{1}
\newcommand\scaleA{0.5}


\section{Results}
\begin{frame}
\frametitle{Without Flow}
\begin{minipage}[t][][t]{0.45\textwidth}
    \begin{figure}
  \centering
\begin{tikzpicture}
  \begin{axis}[
  width=\scalex \textwidth,
  height=\scaley \textwidth,
  xmin=0,
  xmax=1.8,
    ymin=0.4,
    ymax=1,
    ytick distance=0.1,
    grid=major, % Display a grid  
  %grid style={dashed,gray!90}, % Set the style
  xlabel = \small{$ka$},
  ylabel = \small{$|R_{r}|$},
  ]
 \addplot[color=black, thick] table[x index=0,y index=1] {dados/duto_sem_escoamento/analytical_data_abs_r.txt};
 \addplot[color=black, mark=o, only marks] table[x index=0,y index=1] {dados/duto_sem_escoamento/simulation_data_abs_r.txt};
   
  \end{axis}
  \end{tikzpicture}
  %\caption[Coeficiente de Reflexão $R_{r}$ sem Escoamento]{}
  \label{fig:abs_r_sem_escoamento}
\end{figure}
\end{minipage}
\hfill
\begin{minipage}[t]{0.45\textwidth}
\begin{figure}
    \begin{tikzpicture}
  \begin{axis}[
  width=\scalex \textwidth,
  height=\scaley \textwidth,
  xmin=0,
  xmax=1.8,
    ymin=0.4,
    ymax=0.65,
    ytick distance=0.05,
    grid=major, % Display a grid  
  %grid style={dashed,gray!90}, % Set the style
  xlabel = \small{$ka$},
  ylabel = \small{$l/a$},
  ]
 \addplot[color=black, thick] table[x index=0,y index=1] {dados/duto_sem_escoamento/analytical_data_loa.txt};
 \addplot[color=black, mark=o, only marks] table[x index=0,y index=1] {dados/duto_sem_escoamento/simulation_data_loa.txt};
   
  \end{axis}
  \end{tikzpicture}
  %\caption[Coeficiente de Correção da Terminação $l/a$ sem Escoamento]{}
  \label{fig:loa_sem_escoamento}
\end{figure}
\end{minipage}
\centering
Solid lines represents Levine and Schwinger (1948) and circular points are the present study results.
\end{frame}



\begin{frame}
\frametitle{With Outlet Flow: $M = 0.07$ }
\begin{minipage}[t][][t]{0.45\textwidth}
    \begin{figure}
  \centering
\input{figuras/abs_r_exaustao_007.tex}
\end{figure}
\end{minipage}
\hfill
\begin{minipage}[t]{0.45\textwidth}
\begin{figure}
    \begin{tikzpicture}
  \begin{axis}[
   width=\scalex \textwidth,
  height=\scaley \textwidth,
  xmin=0,
  xmax=1.8,
    ymin=0.15,
    ymax=0.6,
    ytick distance=0.05,
    grid=major, % Display a grid  
  %grid style={dashed,gray!90}, % Set the style
  xlabel = \small{$ka$},
  ylabel = \small{$l/a$},
  ]
 \addplot[color=black, thick] table[x index=0,y index=1] {dados/duto_exaustao/loa_007_analytical.txt};
 \addplot[color=black, mark=o, only marks] table[x index=0,y index=1] {dados/duto_exaustao/loa_007_simulation.txt};
   
  \end{axis}
  \end{tikzpicture}
  %\caption[Coeficiente de Correção da Terminação $l/a$ com Escoamento de Exaustão (M $=$ 0,07)]{}
  \label{fig:loa_exaustao_007}
\end{figure}
\end{minipage}
\centering
Solid lines represents Munt (1990) and circular points are the present study results.
\end{frame}

\begin{frame}
\frametitle{With Outlet Flow: $M = 0.15$ }
\begin{minipage}[t][][t]{0.45\textwidth}
    \begin{figure}
  \centering
\input{figuras/abs_r_exaustao_015.tex}
\end{figure}
\end{minipage}
\hfill
\begin{minipage}[t]{0.45\textwidth}
\begin{figure}
    \input{figuras/loa_exaustao_015.tex}
\end{figure}
\end{minipage}
\centering
Solid lines represents Munt (1990) and circular points are the present study results.
\end{frame}

\begin{frame}
\frametitle{With Inlet Flow: $ka = 0.1$}
\begin{minipage}[t][][t]{0.48\textwidth}
    \begin{figure}
  \centering
\begin{tikzpicture}
  \begin{axis}[
  width=\scalex \textwidth,
  height=\scaley \textwidth,
  x tick label style={
      /pgf/number format/.cd,
          fixed,
          fixed zerofill,
          precision=2,
      /tikz/.cd
  },
  xmin=0,
  xmax=0.2,
  ymin=0.68,
  ymax=1,
  ytick distance=0.05,
  xtick distance=0.05,
  grid=major, % Display a grid  
  %grid style={dashed,gray!90}, % Set the style
  xlabel = \small{$M$},
  ylabel = \small{$|R_{r}|$},
  ]
 \addplot[color=black, thick] table[x index=0,y index=1] {dados/duto_sugado/davis_analytical.txt};
 \addplot[color=black, mark=o, only marks] table[x index=0,y index=1] {dados/duto_sugado/davis_simulation.txt};
   
  \end{axis}
  \end{tikzpicture}
  %\caption[Coeficiente de reflexão $R_{M}$ com escoamento sugado]{}

  \label{fig:abs_r_sugado}
\end{figure}
\end{minipage}
\hfill
\begin{minipage}[t]{0.48\textwidth}
\begin{figure}
    \begin{tikzpicture}
  \begin{axis}[
  width=\scalex \textwidth,
  height=\scaley \textwidth,
  x tick label style={
      /pgf/number format/.cd,
          fixed,
          fixed zerofill,
          precision=2,
      /tikz/.cd
  },
  xmin=0,
  xmax=0.2,
  ymin=0.1,
  ymax=1.0,
  ytick distance=0.1,
  xtick distance=0.05,
  grid=major, % Display a grid  
  %grid style={dashed,gray!90}, % Set the style
  xlabel = \small{$M$},
  ylabel = \small{$l/a$},
  ]
 \addplot[color=black, thick] table[x index=0,y index=1] {dados/duto_sugado/davis_analytical_loa.txt};
 \addplot[color=black, mark=o, only marks] table[x index=0,y index=1] {dados/duto_sugado/davis_simulation_loa.txt};
   
  \end{axis}
  \end{tikzpicture}
  %\caption[Coeficiente de correção da terminação ($l_{M}$) com escoamento sugado]{}

  \label{fig:loa_sugado}
\end{figure}
\end{minipage}
\centering
Solid lines represents Davies (1987) and circular points are the present study results.
\end{frame}

\begin{frame}
\frametitle{With Inlet Flow: $K_{p} \text{ X } M$}
\begin{figure}
  \centering
\begin{tikzpicture}
  \begin{axis}[
  width=0.9\textwidth,
  height=0.5\textwidth,
  x tick label style={
      /pgf/number format/.cd,
          fixed,
          fixed zerofill,
          precision=2,
      /tikz/.cd
  },
  xmin=0.05,
  xmax=0.2,
  ymin=0.5,
  ymax=0.77,
  ytick distance=0.05,
  xtick distance=0.03,
  grid=major, % Display a grid  
  %grid style={dashed,gray!90}, % Set the style
  xlabel = $M$,
  ylabel = $K_{p}$,
  ]
 \addplot[color=black, thick] table[x index=0,y index=1] {dados/duto_sugado/kp_m.txt};

  \end{axis}
  \end{tikzpicture}
  %\caption[Coeficiente de reflexão $R_{r}$ com escoamento de exaustão em relação ao número de Mach ($M$) no Strouhal $St = \pi/2$]{Resultado de magnitudes do coeficiente de reflexão $R_{r}$ fixados no Strouhal $St = \pi/2$ em relação ao número de Mach ($M$) para escoamentos sugados. Os resultados foram calculados no ponto $\textbf{P}$ na terminação do duto.}

  \label{fig:kp_m}
\end{figure}
\centering
Loss factor formed by vena contracta increase with Mach number.
\end{frame}

\begin{frame}
\frametitle{With Inlet Flow: $R_{r} \text{ X } ka$}
\begin{figure}
  \centering
\input{figuras/abs_r_sugado_kp.tex}
\end{figure}
\centering
The results with $\bigcirc$ are $M = 0.05$, $\square$ are $M = 0.07$, $\bigtriangleup$ are $M = 0.1$, $\times$ are $M = 0.15$ and $\diamond$ are $M = 0.2$.
\end{frame}

\begin{frame}
\frametitle{With Inlet Flow: $l/a \text{ X } ka$}
\begin{figure}
  \centering
\begin{tikzpicture}
  \begin{axis}[
  width=0.9\textwidth,
  height=0.5\textwidth,
  % x tick label style={
  %     /pgf/number format/.cd,
  %         fixed,
  %         fixed zerofill,
  %         precision=1,
  %     /tikz/.cd
  % },
  xmin=0,
  xmax=1.8,
  ymin=0,
  ymax=1.,
  ytick distance=0.2,
  xtick distance=0.2,
  grid=major, % Display a grid  
  %grid style={dashed,gray!90}, % Set the style
  xlabel = $ka$,
  ylabel = $l/a$,
  ]
 \addplot[color=black, mark=o] table[x index=0,y index=1] {dados/duto_sugado/loa_049_kp.txt};
 \addplot[color=black, mark=square] table[x index=0,y index=1] {dados/duto_sugado/loa_059_kp.txt};
 \addplot[color=black, mark=triangle] table[x index=0,y index=1] {dados/duto_sugado/loa_071_kp.txt};
\addplot[color=black, mark=x] table[x index=0,y index=1] {dados/duto_sugado/loa_074_kp.txt};
\addplot[color=black, mark=diamond] table[x index=0,y index=1] {dados/duto_sugado/loa_075_kp.txt};

  \end{axis}
  \end{tikzpicture}
  %\caption[Coeficiente de reflexão $|R_{r}|$ com escoamento sugado em relação ao número de Strouhal ($St$)]{Resultado de magnitudes do coeficiente de reflexão $|R_{r}|$ em relação ao número de Strouhal ($St$) calculados no ponto $\textbf{P}$ na terminação do duto com vários escoamentos sugados. Calculados pela ferramenta computacional proposta nesse estudo, os pontos com $\bigcirc$ apresentam os resultados para Mach $M =$ 0,05 e número de Reynolds $Re =$ 1378,73, os pontos com $\square$ apresentam os resultados para Mach $M =$ 0,07 e número de Reynolds $Re =$ 1930,23, os pontos com $\bigtriangleup$ apresentam os resultados para Mach $M =$ 0,1 e número de Reynolds $Re =$ 2757,42, os pontos com $\times$ apresentam os resultados para Mach $M =$ 0,15 e número de Reynolds $Re =$ 2057,71 e os pontos com $\diamond$ apresentam os resultados para Mach $M =$ 0,20 e número de Reynolds $Re =$ 5514,82.}

  \label{fig:loa_sugado_kp}
\end{figure}
\centering
The results with $\bigcirc$ are $M = 0.05$, $\square$ are $M = 0.07$, $\bigtriangleup$ are $M = 0.1$, $\times$ are $M = 0.15$ and $\diamond$ are $M = 0.2$.
\end{frame}

\section{Conclusions}
\begin{frame}
\frametitle{Conclusions}
\begin{itemize}
  \item Without flow: Good agreements were obtained among results from model in lattice Boltzmann method and Levine and Schwinger (1948);
  \item Outlet flow: Good agreements were obtained among results from model in lattice Boltzmann method and Munt (1990);
  \item Inlet flow: Good agreements were obtained among $|R{r}|$ results from model in lattice Boltzmann method and Davies (1997). But there are divergences with values of $l/a$;
  \item Vena contracta was analyzed and your factor loss $K_{p}$ increase with inlet flow Mach;
  \item In general $|R_{r}|$ decrease with increased values of $K_{p}$;
  \item In general $l/a$ don't change with differents values of $K_{p}$;
  \item By increase values of $M$, $|R_{r}|$ increase with outlet flow and decrease with intlet flow.
\end{itemize}
\end{frame}

%-------------------------------------

\begin{frame}
\begin{center}
Thank you!\\[2cm]Questions?
\end{center}
\end{frame}

\end{document}


